% Ad M. Brutum 2.4 --- headnote
% ad-brutum-02-04, 12 April 43 BC, Rome

\textit{Cicero to M. Junius Brutus, written from Rome at
dawn --- Perseus dateline \textit{Scr. Romae prid. Id.
Apr. mane a. 711 (43)}, i.e.\ 12 April 43 BC, in the
morning. The dateline matches the meta date. The letter
is an explicit postscript: the previous day's letter (to
which 2.4 is functionally a sequel) was given to Scaptius
for delivery on 11 April, and a letter from Brutus,
written at Dyrrachium on 1 April, was delivered to
Cicero that same evening; when Cicero learned next
morning that Scaptius's messengers had not yet
departed, he scratched out this supplement \textit{in
ipsa turba matutinae salutationis} --- ``amid the very
throng of my morning callers'' --- a phrase that
preserves the texture of the consular morning. The
letter is therefore a snapshot of a single Roman dawn
in the worst week of the Mutina campaign: two days
before Forum Gallorum, written before the householder
has finished receiving his clients.}

\textit{The letter answers Brutus point by point. Section
2 reports the Senate's decree authorising Cassius to
pursue Dolabella, won over the angry opposition of the
consul Pansa --- a notable detail, since Pansa is
within a fortnight of dying of his Forum-Gallorum wounds.
Section 3 is the most politically pointed of the
six paragraphs: Brutus has asked what to do about his
prisoner C.\ Antonius (Mark Antony's brother, captured
at Apollonia), and Cicero replies, with the bare
sufficiency he reserves for matters where Brutus
already knows his mind, ``until we have learned how
Brutus's case has come out, I think he should be kept in
custody.'' The Perseus text in this section carries an
unresolved crux at \textit{eam semel cepit, mihi crede,
non erit id.\ Apr.}, here preserved between daggers.
Section 4 is the practical answer about troops and
money: Brutus must borrow from the cities, and Pansa
will not release reinforcements --- whether out of
strategic prudence or out of jealousy of Brutus's
strength, Cicero leaves open. Section 5 deals with the
spread of Cassius's news in Rome and the temper of
the still-existing \textit{partes Caesaris}; and section
6 is the famous coda about Cicero's son, then on
Brutus's staff in Macedonia, in which the father's
delight at the son's promise is openly inseparable
from his delight that Brutus loves him. The letter is
one of the most domestically intimate of the Brutus
correspondence and one of the most politically dense.}
